\section*{Bab \ref{ch:b1:taylor} --- Rumus Taylor dan Ekspansi
Asimtotik}

\begin{solution}{exo:b1:taylor:1}
\[
\eu^{2x} = 1 + 2x + 2x^2 + \frac{4x^3}{3} + o(x^3);
\qquad
\ln(1 - x) = -x - \frac{x^2}{2} - \frac{x^3}{3} - \frac{x^4}{4} +
o(x^4);
\]
\[
\sqrt{1+x} = 1 + \frac x2 - \frac{x^2}{8} + \frac{x^3}{16} + o(x^3);
\qquad
\frac{1}{1 + x^2} = 1 - x^2 + x^4 - x^6 + o(x^6),
\]
yang terakhir lewat penyulihan $-x^2$ ke dalam ekspansi geometrinya.
\end{solution}

\begin{solution}{exo:b1:taylor:2}
$\eu^x - 1 - x = \frac{x^2}{2} + o(x^2)$: jadi limitnya $\dfrac12$.

$\cos x = 1 - \frac{x^2}{2} + \frac{x^4}{24} + o(x^4)$ dan
$\sqrt{1 - x^2} = 1 - \frac{x^2}{2} - \frac{x^4}{8} + o(x^4)$:
sehingga selisihnya $\frac{x^4}{24} + \frac{x^4}{8} = \frac{x^4}{6} + o(x^4)$:
jadi limitnya $\dfrac16$.

$\ln(1+x) - \sin x = \bigl(x - \frac{x^2}{2}\bigr) - x + o(x^2) =
-\frac{x^2}{2} + o(x^2)$: jadi limitnya $-\dfrac12$.
\end{solution}

\begin{solution}{exo:b1:taylor:3}
$\frac{1}{1+t^2} = 1 - t^2 + t^4 + o(t^5)$; lalu mengintegralkannya dari $0$ sampai
$x$ (\cref{met:b1:taylor:compute}~(5)):
\[
\arctan x = x - \frac{x^3}{3} + \frac{x^5}{5} + o(x^5)\ \ (\text{bahkan
}o(x^6)\text{, karena ganjil}).
\]
$(1 - t^2)^{-1/2} = 1 + \frac{t^2}{2} + \frac38 t^4 + o(t^5)$ (lewat ekspansi
binomial dengan $\alpha = -\frac12$, $x = -t^2$:
$\binom{-1/2}{2} = \frac{(-\frac12)(-\frac32)}{2} = \frac38$);
lalu mengintegralkannya:
\[
\arcsin x = x + \frac{x^3}{6} + \frac{3x^5}{40} + o(x^5) .
\]
\end{solution}

\begin{solution}{exo:b1:taylor:4}
Taylor--Lagrange (\cref{thm:b1:taylor:lagrange}) bagi $\exp$ di $a =
0$, $x = 1$: \omterm{def:b1:derivative:def}{turunan} ke-$(n+1)$-nya adalah $\eu^t \leq \eu < 3$ pada
$\intcc{0}{1}$, sehingga
\[
\Bigl|\eu - \sum_{k=0}^{n} \frac{1}{k!}\Bigr| \leq
\frac{3}{(n+1)!} .
\]
Untuk $6$ desimal yang persis, kita menginginkan $\frac{3}{(n+1)!} < 5\times 10^{-7}$,
yakni $(n+1)! > 6\times 10^{6}$: dan karena $10! = 3\,628\,800$ dan $11! =
39\,916\,800$, maka $n + 1 = 11$, yakni $n = 10$ mencukupi.
\end{solution}

\begin{solution}{exo:b1:taylor:5}
$n\ln\bigl(1 + \frac1n\bigr) = n\Bigl(\frac1n - \frac{1}{2n^2} +
\frac{1}{3n^3} + o\bigl(\frac{1}{n^3}\bigr)\Bigr) = 1 - \frac{1}{2n}
+ \frac{1}{3n^2} + o\bigl(\frac{1}{n^2}\bigr)$. Lalu mengeksponensialkannya, dengan
$u = -\frac{1}{2n} + \frac{1}{3n^2}$ dan $\eu^u = 1 + u + \frac{u^2}2
+ o(u^2)$:
\[
\Bigl(1 + \frac1n\Bigr)^n = \eu\cdot \eu^{u}
= \eu\Bigl(1 - \frac{1}{2n} + \frac{1}{3n^2} + \frac{1}{8n^2}
+ o\Bigl(\frac{1}{n^2}\Bigr)\Bigr)
= \eu\Bigl(1 - \frac{1}{2n} + \frac{11}{24n^2} +
o\Bigl(\frac{1}{n^2}\Bigr)\Bigr).
\]
Jadi limitnya $\eu$; sedangkan galatnya $\sim \dfrac{\eu}{2n}$: yang lamban (yakni satu angka per
kenaikan sepuluh kali lipat $n$).
\end{solution}

\begin{solution}{exo:b1:taylor:6}
$f(x) = x^2 - x^4 = x^2(1 + o(1))$: dengan suku pertamanya $x^2$, $p = 2$ yang genap,
dan koefisiennya $> 0$: jadi minimum lokal di $0$ (yang tak menyeluruh: karena $f(2) = -12$).

$g(x) = x^3 + x^5$: dengan suku pertamanya $x^3$, dan $p$ yang ganjil: jadi tak ada ekstremumnya; sehingga $g$
memotong garis singgungnya (yang mendatar): yakni titik belok di $0$.

$h = \exp$ di $a = 1$: $h(x) = \eu + \eu(x-1) + \frac{\eu}{2}(x-1)^2
+ o((x-1)^2)$; sehingga selisihnya dengan garis singgungnya adalah
$\frac{\eu}{2}(x-1)^2 + o(\cdot) > 0$ secara setempat. Adapun secara menyeluruh: $\eu^x -
\eu x \geq 0$ untuk setiap $x$ menurut kecembungannya
(\cref{thm:b1:derivative:convexchar}~(3)): jadi grafiknya terletak di atas setiap
garis singgungnya, dengan kesamaannya hanya di titik sentuhnya.
\end{solution}

\begin{solution}{exo:b1:taylor:7}
Untuk $x \to +\infty$:
\[
f(x) = x\Bigl(1 + \frac1x\Bigr)^{1/3}
= x\Bigl(1 + \frac{1}{3x} - \frac{1}{9x^2} +
o\Bigl(\frac{1}{x^2}\Bigr)\Bigr)
= x + \frac13 - \frac{1}{9x} + o\Bigl(\frac1x\Bigr):
\]
jadi asimtotnya $y = x + \frac13$, dengan kurvanya di bawahnya di dekat $+\infty$. Adapun ketika $x
\to -\infty$, perhitungan yang sama tetap sah (karena akar pangkat tiganya
terdefinisi bagi setiap bilangan real, dan $\frac1x \to 0$): jadi asimtot yang sama $y = x +
\frac13$, tetapi kini $-\frac{1}{9x} > 0$: sehingga kurvanya \emph{di atas} garisnya.
\end{solution}

\begin{solution}{exo:b1:taylor:8}
$u_n = \sqrt n\bigl(\sqrt{1 + \tfrac1n} - 1\bigr) = \sqrt
n\bigl(\frac{1}{2n} + o(\frac1n)\bigr) \sim \dfrac{1}{2\sqrt n}$.

$v_n = \ln\bigl(1 + \frac1n\bigr) \sim \dfrac 1n$.

Adapun $w_n$: dengan $h = \frac1n \to 0$, $\sin h - \tan h = \bigl(h -
\frac{h^3}{6}\bigr) - \bigl(h + \frac{h^3}{3}\bigr) + o(h^3) =
-\frac{h^3}{2} + o(h^3)$, sehingga $w_n \sim -\dfrac{1}{2n^3}$.
\end{solution}

\begin{solution}{exo:b1:taylor:9}
Taylor--Lagrange pada orde $1$ di sekitar $x$, pada kedua sisinya:
\[
f(x + h) = f(x) + h f'(x) + R_+,\quad
f(x - h) = f(x) - h f'(x) + R_-,
\qquad \abs{R_\pm} \leq \frac{h^2}{2} \sup \abs{f''} .
\]
Lalu menguranginya: $f(x+h) - f(x-h) = 2h f'(x) + (R_+ - R_-)$, sehingga
\[
\abs{f'(x)} \leq \frac{\abs{f(x+h) - f(x-h)}}{2h} +
\frac{h}{2}\sup_{\intcc{x-h}{x+h}}\abs{f''} .
\]
Dengan batas yang menyeluruh: $\abs{f'(x)} \leq \frac{M_0}{h} + \frac{M_2
h}{2}$ untuk setiap $h > 0$. Adapun ruas kanannya minimum di $h =
\sqrt{2M_0/M_2}$ (karena \omterm{def:b1:derivative:def}{turunannya} nol), dengan nilai $\sqrt{2M_0M_2}$
--- sehingga $\abs{f'} \leq \sqrt{2M_0M_2}$. (Jika $M_2 = 0$, biarkanlah $h \to
\infty$: maka $f' = 0$, yang selaras.)
\end{solution}

\begin{solution}{exo:b1:taylor:10}
Pada $\intoo{0}{\pi}$: $0 < \sin u < u$, sehingga $(u_n)$ turun
tegas, positif, jadi konvergen; dan limitnya merupakan titik tetap
$\sin$ di $\intcc{0}{\pi}$, sedangkan $\sin \ell = \ell$ memaksa $\ell =
0$ (karena $\sin x < x$ untuk $x > 0$).
\begin{enumerate}
  \item Dengan memakai $\sin u = u - \frac{u^3}{6} + o(u^3)$ ketika $u \to 0$:
        \[
        v_{n+1} - v_n = \frac{1}{\sin^2 u_n} - \frac{1}{u_n^2}
        = \frac{1}{u_n^2}\Bigl(\Bigl(1 - \frac{u_n^2}{6} +
        o(u_n^2)\Bigr)^{\!-2} - 1\Bigr)
        = \frac{1}{u_n^2}\Bigl(\frac{u_n^2}{3} + o(u_n^2)\Bigr)
        \longrightarrow \frac13 .
        \]
  \item Menurut \cref{exo:b1:seq:10}~(3) (yaitu Cesàro bagi selisihnya),
        $\frac{v_n}{n} \to \frac13$, yakni $v_n \sim \frac n3$, yakni
        $u_n^2 \sim \frac 3n$: dan karena $u_n > 0$,
        \[
        u_n \sim \sqrt{\frac{3}{n}} .
        \]
\end{enumerate}
\end{solution}

\begin{solution}{exo:b1:taylor:11}
Lewat penyebut bersamanya. Limit pertamanya:
\[
\frac{1}{x^2} - \frac{1}{\sin^2 x}
= \frac{\sin^2 x - x^2}{x^2\sin^2 x},
\qquad
\sin^2 x = \Bigl(x - \frac{x^3}{6} + o(x^4)\Bigr)^{\!2}
= x^2 - \frac{x^4}{3} + o(x^5) ,
\]
sehingga pembilangnya $-\frac{x^4}{3} + o(x^4)$ sedangkan
penyebutnya $\sim x^4$: jadi limitnya $-\dfrac13$.

Yang kedua: $\dfrac1x - \dfrac{1}{\ln(1+x)} = \dfrac{\ln(1+x) -
x}{x\ln(1+x)}$; dengan pembilangnya $-\frac{x^2}{2} + o(x^2)$, dan
penyebutnya $x\bigl(x + o(x)\bigr) \sim x^2$: jadi limitnya
$-\dfrac12$.
\end{solution}

\begin{solution}{exo:b1:taylor:12}
Pada $I_k = \intoo{k\pi - \frac\pi2}{k\pi + \frac\pi2}$, fungsi
$g(x) = \tan x - x$ berturunan $\tan^2 x \geq 0$,
yang lenyap hanya di titik tunggal $k\pi$: sehingga $g$ naik
tegas pada $I_k$ (\cref{cor:b1:derivative:monotone}~(2)),
dengan limit $-\infty$ dan $+\infty$ di ujungnya: jadi tepat satu
nol $x_k$. Untuk $k \geq 1$, $g(k\pi) = -k\pi < 0$, sehingga $x_k \in
\intoo{k\pi}{k\pi + \frac\pi2}$: tulislah $x_k = k\pi + \frac\pi2 -
\varepsilon_k$ dengan $\varepsilon_k \in \intoo{0}{\frac\pi2}$.
Lalu
\[
x_k = \tan x_k = \tan\Bigl(\frac\pi2 - \varepsilon_k\Bigr)
= \frac{1}{\tan\varepsilon_k}
\quad\Longrightarrow\quad
\varepsilon_k = \arctan\frac{1}{x_k} ,
\]
dengan memakai $\tan\varepsilon_k = \frac{1}{x_k}$ dan $\varepsilon_k \in
\intoo{0}{\frac\pi2}$. Dan karena $x_k \geq k\pi \to \infty$:
maka $\varepsilon_k \to 0$, sehingga
\[
\varepsilon_k = \arctan\frac{1}{x_k} = \frac{1}{x_k} +
O\Bigl(\frac{1}{x_k^3}\Bigr)
= \frac{1}{k\pi + O(1)} + O\Bigl(\frac{1}{k^3}\Bigr)
= \frac{1}{k\pi} + O\Bigl(\frac{1}{k^2}\Bigr) ,
\]
sehingga $x_k = k\pi + \frac\pi2 - \frac{1}{k\pi} +
o\bigl(\frac1k\bigr)$.
\end{solution}

\begin{solution}{pb:b1:taylor:1}
\textbf{1.} Di sini $S_{2n+1} - S_{2n-1} = a_{2n} - a_{2n+1} \geq 0$ dan
$S_{2n+2} - S_{2n} = a_{2n+2} - a_{2n+1} \leq 0$, sedangkan $S_{2n}
- S_{2n+1} = a_{2n+1} \to 0$: sehingga barisan $(S_{2n+1})$ dan
$(S_{2n})$ berdampingan, yang konvergen ke $S$ yang sama
(\cref{thm:b1:seq:adjacent}) dengan $S_{2n+1} \leq S \leq S_{2n}$.
Untuk $n$ yang genap: $S_{n+1} \leq S \leq S_n$ memberikan $-a_{n+1} \leq S -
S_n \leq 0$; sedangkan untuk $n$ yang ganjil: $0 \leq S - S_n \leq a_{n+1}$. Pada kedua
kasusnya $\abs{S - S_n} \leq a_{n+1}$ dan $S - S_n$ bertanda sama dengan
$(-1)^{n+1}$, yaitu suku pertama yang dihilangkan. Adapun penurunan yang tegas membuat
setiap ketaksamaan yang terpampang menjadi tegas, khususnya $0 < \abs{S -
S_n} < a_{n+1}$.

\textbf{2.} Di sini $a_k = \frac{1}{k!}$ turun tegas ke $0$:
sehingga pertanyaan 1 berlaku. Lalu Taylor--Lagrange
(\cref{thm:b1:taylor:lagrange}) bagi $\exp$ di antara $-1$ dan $0$:
$\abs{\eu^{-1} - T_n} \leq \frac{1}{(n+1)!}$ (karena \omterm{def:b1:derivative:def}{turunannya}
$\eu^t$ bernilai $\leq 1$ di sana), sehingga $T_n \to \eu^{-1}$, dan limit
$S$ pada pertanyaan 1 \emph{adalah} $\eu^{-1}$, dengan batas tegasnya
$0 < \abs{\eu^{-1} - T_n} < \frac{1}{(n+1)!}$.

\textbf{3.} Dengan mengintegralkan kesamaannya atas $\intcc{0}{1}$: ruas
kirinya menjadi $\arctan 1 = \frac\pi4$ (menurut teorema fundamentalnya), lalu
suku ke-$k$-nya memberikan $\frac{(-1)^k}{2k+1}$, dan
\[
\rho_n = (-1)^{n+1}\int_0^1 \frac{t^{2n+2}}{1+t^2}\dd t,
\qquad
\abs{\rho_n} \leq \int_0^1 t^{2n+2}\dd t = \frac{1}{2n+3} .
\]

\textbf{4.} Galat pada $\pi$ adalah $4\abs{\rho_n} \leq
\frac{4}{2n+3}$: dan di bawah $10^{-6}$ menuntut $2n + 3 > 4\cdot10^6$,
yakni kira-kira dua juta suku. Sementara itu $4S_4 = 4\bigl(1 - \frac13 +
\frac15 - \frac17 + \frac19\bigr) = 4 \times 0.834921 =
3.339683$, yang hampir $0.2$ jauhnya dari $\pi$: jadi lima suku, tanpa
satu angka pun.

\textbf{5.} Dengan mengintegralkannya atas $\intcc{0}{1}$: $\ln 2 =
\sum_{k=1}^{n} \frac{(-1)^{k-1}}{k} + (-1)^n R_n$ dengan $R_n =
\int_0^1 \frac{t^n}{1+t}\dd t$; lalu dari $\frac12 \leq \frac{1}{1+t}
\leq 1$: $\frac{1}{2(n+1)} \leq R_n \leq \frac{1}{n+1}$. Jadi
galatnya terperangkap di antara dua kelipatan $\frac1n$: sehingga enam
desimal berongkos kira-kira sejuta suku.

\textbf{6.} Dengan mengintegralkannya dari $0$ sampai $x$: $\frac12\ln\frac{1 +
x}{1 - x} = \sum_{k=0}^{n} \frac{x^{2k+1}}{2k+1} + \int_0^x
\frac{t^{2n+2}}{1-t^2}\dd t$. Lalu di $x = \frac13$: $\frac{1 +
1/3}{1 - 1/3} = 2$, dan pada $\intcc{0}{\frac13}$, $\frac{1}{1 -
t^2} \leq \frac98$:
\[
\ln 2 = 2\sum_{k=0}^{n} \frac{(1/3)^{2k+1}}{2k+1} +
\tilde\rho_n,
\qquad
0 < \tilde\rho_n \leq \frac{9}{4}\cdot
\frac{(1/3)^{2n+3}}{2n+3} :
\]
sehingga setiap suku tambahannya \omterm{def:b1:arith:divides}{membagi} galatnya kira-kira dengan $9$.

\textbf{7.} Lewat induksi: untuk $n = 1$: $1 - \frac12 = \frac12 = H_2
- H_1$. Adapun langkahnya:
\[
H_{2n+2} - H_{n+1} = (H_{2n} - H_n) + \frac{1}{2n+1} +
\frac{1}{2n+2} - \frac{1}{n+1}
= (H_{2n} - H_n) + \frac{1}{2n+1} - \frac{1}{2n+2},
\]
yang persis merupakan pertambahan jumlah berselang-selingnya. Dan
$H_{2n} - H_n = \sum_{k=1}^{n}\frac{1}{n+k}$ merupakan \omterm{thm:b1:integration:riemann}{jumlah Riemann}
pada \cref{ex:b1:integration:riemannexample}, yang konvergen ke $\ln
2$: jadi jumlah parsial genap Cara 1 \emph{adalah} \omterm{thm:b1:integration:riemann}{jumlah Riemann}
Cara 3.

\textbf{8.} Di sini $\sum_{k=1}^{6}\frac{(-1)^{k-1}}{k} = 0.61667$,
dengan galat $0.0765$; sedangkan Cara 2 pada $n = 5$ memberikan $0.6931471$ dengan galat
$\leq \frac94\cdot\frac{(1/3)^{13}}{13} = 1.1\cdot10^{-7}$. Adapun
alasannya: bahwa Cara 1 menilai deret logaritmanya di titik
\omterm{def:b1:topology:closure}{perbatasan} $x = 1$, yang di situ sukunya meluruh seperti $\frac1k$; sedangkan Cara 2
menilainya di $x = \frac13$, yang jauh di dalam, yang di situ setiap sukunya membawa
faktor $\frac19$ yang baru.

\textbf{9.} Sepuluh desimal: kita menginginkan $\tilde\rho_n \leq
5\cdot10^{-11}$. Pada $n = 10$: $\frac94 \cdot
\frac{(1/3)^{23}}{23} = \frac94\cdot\frac{1.06\cdot10^{-11}}{23}
\approx 1.0\cdot10^{-12} < 5\cdot10^{-11}$: sehingga sebelas suku
mencukupi.

\textbf{10.} Di sini $(5+\iu)^2 = 24 + 10\iu$, lalu $(5+\iu)^4 = (24 +
10\iu)^2 = 476 + 480\iu$; dan $2(1+\iu)(239+\iu) = 2(238 +
240\iu) = 476 + 480\iu$: jadi sama. Adapun argumennya: $\arg(5 + \iu) =
\arctan\frac15$, sehingga ruas kirinya berargumen
$4\arctan\frac15 \approx 0.79 \in \intoo{0}{\pi}$; sedangkan ruas
kanannya berargumen $\frac\pi4 + \arctan\frac{1}{239} \in
\intoo{0}{\pi}$. Jadi dua bilangan kompleks yang sama dengan argumen pada
\omterm{prop:b1:reals:intervals}{selang} yang sama berpanjang $< 2\pi$:
\[
4\arctan\frac15 = \frac\pi4 + \arctan\frac{1}{239} ,
\]
yang merupakan rumus Machin.

\textbf{11.} Integralkanlah $\frac{1}{1+t^2} = \sum_{k=0}^n (-1)^k
t^{2k} + \frac{(-1)^{n+1}t^{2n+2}}{1+t^2}$ dari $0$ sampai $x$:
\[
\arctan x = \sum_{k=0}^{n}\frac{(-1)^k x^{2k+1}}{2k+1} +
r_n(x),
\qquad
\abs{r_n(x)} \leq \int_0^x t^{2n+2}\dd t =
\frac{x^{2n+3}}{2n+3} .
\]

\textbf{12.} Galatnya: $16\,\frac{(1/5)^{11}}{11} = 3.0\cdot
10^{-8}$ dan $4\,\frac{(1/239)^5}{5} \approx 1\cdot10^{-12}$:
jadi totalnya $< 5\cdot10^{-8}$. Adapun jumlah yang terpampang bernilai
$3.14159268\dots$, sehingga $\pi = 3.1415926\dots$ tersertifikasi sampai
$5\cdot10^{-8}$: jadi tujuh desimal dari enam suku (yaitu lima di
$\frac15$, dan dua di $\frac1{239}$ bila dihitung dengan murah hati).

\textbf{13.} Leibniz: dengan galat $\sim \frac1n$, sehingga setiap angka barunya
mengalikan bebannya dengan sepuluh. Dalzell
(\cref{pb:b1:integration:1}, pertanyaan 22): dengan galat $4^{1-5m}$,
yakni kira-kira tiga angka per langkah, dengan setiap langkahnya berupa \omterm{def:b1:poly:def}{polinomial} yang lebih berat.
Machin: dengan nisbah galat $\frac{1}{25}$ per suku, yakni kira-kira $1.4$ angka
per suku, dengan setiap sukunya satu pembagian. Adapun moralnya: bahwa sisa
ekspansi bertipe geometri terskala seperti $x^{2n}$, sehingga membuat $x$
kecil membeli angka dengan ongkos yang \emph{tetap} per suku --- jadi kesamaan
kompleks Machin persis merupakan mesin bagi pengecilan $x$.

\textbf{14.} Di sini $a_k = \frac{1}{(2k)!}$ turun tegas ke $0$;
sehingga menurut pertanyaan 1 dan \cref{prop:b1:taylor:standard} (dengan batas
Lagrange seperti pada pertanyaan 2), $\sum_{k \leq n}\frac{(-1)^k}{(2k)!}
\to \cos 1$ dengan kurung yang \emph{tegas} $0 < \bigl|\cos 1 -
S'_n\bigr| < \frac{1}{(2n+2)!}$. Andaikan $\cos 1 = \frac pq$ lalu
ambillah $2n \geq q$: maka $(2n)!\,S'_n = \sum_{k\leq n} (-1)^k
\frac{(2n)!}{(2k)!} \in \Z$ dan $(2n)!\,\frac pq \in \Z$, sedangkan
\[
0 < \Bigl|(2n)!\cos 1 - (2n)!S'_n\Bigr| <
\frac{(2n)!}{(2n+2)!} = \frac{1}{(2n+1)(2n+2)} < 1 :
\]
yaitu bilangan bulat taknol yang nilai mutlaknya $< 1$. Jadi bertentangan: sehingga $\cos
1 \notin \Q$.

\textbf{15.} Dengan cara yang identik memakai $a_k = \frac{1}{(2k+1)!}$,
lalu mengalikannya dengan $(2n+1)!$ dengan $2n + 1 \geq q$: maka $\sin 1 \notin
\Q$. Namun $\cos^2 1 + \sin^2 1 = 1 \in \Q$: jadi hasil kali dan jumlah
bilangan irasional boleh saja rasional --- sehingga keirasionalan tak melewati
operasi aljabar mana pun secara cuma-cuma.

\textbf{16.} Adapun batas ekornya: untuk $m > n$,
\[
\sum_{k=n+1}^{m} \frac{1}{(2k)!}
\leq \frac{1}{(2n+2)!}\Bigl(1 + \frac12 + \frac14 +
\dots\Bigr) \leq \frac{2}{(2n+2)!} ,
\]
karena setiap nisbah berurutannya adalah $\frac{1}{(2k+1)(2k+2)} \leq
\frac12$; sedangkan ekornya positif (karena suku pertamanya demikian). Jadi $0 <
\cosh 1 - \sum_{k\leq n}\frac{1}{(2k)!} < \frac{2}{(2n+2)!}$,
lalu mengalikannya dengan $(2n)!$ dengan $2n \geq q$ memerangkap bilangan bulat
taknol di $\intoo{0}{1}$ lagi: sehingga $\cosh 1 \notin \Q$.

\textbf{17.} Di sini $\cos\frac1m = \sum_k
\frac{(-1)^k}{m^{2k}(2k)!}$: sukunya turun tegas ke
nol, dan $m^{2n}(2n)!\cdot\frac{1}{m^{2k}(2k)!} =
m^{2(n-k)}\frac{(2n)!}{(2k)!} \in \Z$ untuk $k \leq n$. Jika
$\cos\frac1m = \frac pq$, kalikanlah kurung tegasnya dengan
$q\,m^{2n}(2n)!$: maka galatnya terbatas oleh
$\frac{q}{m^2(2n+1)(2n+2)} < 1$ untuk $n$ yang besar: jadi bertentangan.
Adapun untuk $\frac ab$ dengan $a \geq 2$: membersihkan penyebutnya mengalikan
ekornya dengan $b^{2n}(2n)!$, tetapi suku pertama yang dihilangkan adalah
$\frac{a^{2n+2}}{b^{2n+2}(2n+2)!}$, dan hasil kalinya
$\frac{a^{2n+2}}{b^2(2n+1)(2n+2)}$ \emph{meledak}: sehingga
pembilangnya $a^{2n+2}$ tak lagi terbersihkan, dan perangkapnya macet.
(Adapun hasilnya tetap benar --- lewat mesin bergaya Niven, bukan
yang ini.)

\textbf{18.} Di $t = x$ setiap suku $g$ lenyap kecuali $f(x)
- f(x) = 0$: sehingga $g(x) = 0$; dan $A$ dipilih agar $g(a) = 0$.
Lalu menurunkannya, jumlahnya berteleskop:
\[
g'(t) = -\frac{f^{(n+1)}(t)}{n!}(x - t)^n +
A\,\frac{(x-t)^n}{n!}
= \frac{(x-t)^n}{n!}\bigl(A - f^{(n+1)}(t)\bigr) .
\]
Lalu Rolle pada ruas dari $a$ ke $x$ memberikan $c$ yang tegas di antaranya
dengan $g'(c) = 0$; dan karena $(x - c)^n \neq 0$: maka $A =
f^{(n+1)}(c)$. Sehingga membuka $g(a) = 0$ menghasilkan kesamaan Taylornya
dengan sisa $\frac{f^{(n+1)}(c)}{(n+1)!}(x-a)^{n+1}$.

\textbf{19.} Untuk $x > 0$, sisanya adalah $\frac{\eu^{c}}
{(n+1)!}x^{n+1} > 0$: sehingga eksponensialnya melampaui setiap
\omterm{def:b1:poly:def}{polinomial} Taylornya, secara tegas, pada setiap ordenya. Adapun untuk $x < 0$
tanda sisanya sama dengan tanda $x^{n+1}$: jadi $\eu^x$ berada
\emph{di atas} \omterm{def:b1:poly:def}{polinomialnya} untuk $n$ yang ganjil, dan \emph{di bawahnya} untuk
$n$ yang genap --- yaitu sisi yang berselang-seling, sebagaimana grafik $1 + x$ dan $1
+ x + \frac{x^2}{2}$ terhadap $\eu^x$ sudah memperlihatkan.

\textbf{20.} Galat sejatinya: $\sin 0.5 - 0.4791667 =
2.59\cdot10^{-4}$, terhadap batasnya $\frac{0.5^5}{120} =
2.60\cdot10^{-4}$: jadi hampir tercapai (karena suku berikutnya mendominasi
ekornya). Adapun Young: bagi limit dan analisis setempat, yang di situ tak ada konstanta
yang diperlukan. Lagrange: bagi desimal yang bersertifikat. Sedangkan yang berselang-seling: bila
berlaku, ia memberikan batas yang sama \emph{ditambah} arah galatnya
--- jadi yang terbaik di antara ketiganya, tetapi yang paling jarang.

\textbf{21.} \omterm{def:b1:continuity:continuous}{Kekontinuannya} di $0$: dengan $u = \frac{1}{x^2} \to
+\infty$, $f(x) = \eu^{-u} \to 0 = f(0)$. Adapun \omterm{def:b1:derivative:def}{turunannya} di $0$:
$\bigl|\frac{f(h)}{h}\bigr| = \sqrt u\,\eu^{-u} \to 0$
(\cref{prop:b1:functions:powerrules}): sehingga $f'(0) = 0$. Untuk $x \neq
0$, $f'(x) = \frac{2}{x^3}\eu^{-1/x^2}$: yaitu bentuk
$P_1\bigl(\frac1x\bigr)\eu^{-1/x^2}$ dengan $P_1(X) = 2X^3$;
lalu secara induktif, menurunkan $P_k(\frac1x)\eu^{-1/x^2}$ memberikan
$P_{k+1}(X) = 2X^3 P_k(X) - X^2 P_k'(X)$, yaitu sebuah \omterm{def:b1:poly:def}{polinomial}. Lalu
\[
\frac{f^{(k)}(h) - 0}{h} = v\,P_k(v)\,\eu^{-v^2}
\Big|_{v = 1/h} \longrightarrow 0
\]
(yaitu \omterm{def:b1:poly:def}{polinomial} terhadap $\eu^{-v^2}$, menurut \omterm{prop:b1:functions:powerrules}{perbandingan pertumbuhan} di
$\pm\infty$): sehingga lewat induksi $f^{(k)}(0) = 0$ untuk setiap $k$. Jadi semua
\omterm{def:b1:poly:def}{polinomial} Taylor $f$ di $0$ lenyap, namun $f > 0$ di luar $0$:
sehingga Taylor--Young persis pada setiap ordenya dan buta di luar
perilaku setempatnya. Jadi sebuah ekspansi hanyalah keterangan yang setempat.

\textbf{22.} Menurut pertanyaan 2, $0 < \abs{\eu^{-1} - T_n} <
\frac{1}{(n+1)!}$, secara tegas. Jika $\eu^{-1} = \frac pq$, ambillah $n
\geq q$ lalu kalikanlah dengan $n!$: maka $n!\,T_n \in \Z$ dan $n!\frac pq
\in \Z$, sehingga sebuah bilangan bulat taknol bernilai mutlak $<
\frac{n!}{(n+1)!} = \frac{1}{n+1} < 1$: jadi bertentangan. Sehingga
$\eu^{-1} \notin \Q$, dan $\eu = \frac{1}{\eu^{-1}}$ pun
irasional. Adapun ketiga buktinya: \omterm{thm:b1:seq:adjacent}{barisan berdampingan} yang mengapit
$q!\,\eu$ (\cref{exo:b1:seq:9}); rekurensi \omterm{thm:b1:integration:def}{integral} $A_n =
\eu - nA_{n-1}$ (\cref{pb:b1:integration:1}); dan kurung
berselang-selingnya (di sini). Jadi satu perangkap, tiga sertifikat kekecilan.

\textbf{23.} Kelompokkanlah jumlah parsialnya berpasangan: $\sum_{k=1}^{2n}
(-1)^k b_k = E_n - O_n$ dengan $E_n = \sum_{j=1}^{n}
\frac{1}{4j^2}$, yang terbatas (menurut batas teleskopis pada
\cref{ex:b1:seq:basel}, $E_n \leq \frac12$), dan $O_n =
\sum_{j=1}^{n}\frac{1}{2j-1} \geq \frac12 H_n \to +\infty$
(\cref{exo:b1:seq:5}): sehingga jumlah parsialnya menuju $-\infty$.
Adapun kemonotonannya dipakai pada pertanyaan 1 persis di tempat $S_{2n+1} -
S_{2n-1} = a_{2n} - a_{2n+1}$ memerlukan sebuah tanda: karena tanpa penurunannya,
\omterm{def:b1:seq:subsequence}{subbarisan} genap dan ganjilnya tak harus monoton, sehingga
kedampingannya runtuh.

\textbf{24.} Di sini $(2+\iu)(3+\iu) = 5 + 5\iu = 5(1+\iu)$; lalu mengambil
argumennya (yang semuanya di $\intoo{0}{\frac\pi2}$): $\arctan\frac12 +
\arctan\frac13 = \frac\pi4$. Adapun ongkos deretnya bagi enam desimal:
galatnya $\leq 4\bigl(\frac{(1/2)^{2n+3}}{2n+3} +
\frac{(1/3)^{2n+3}}{2n+3}\bigr)$, yang pada $n = 10$ bernilai
$\approx 2\cdot10^{-8} < 5\cdot10^{-7}$: jadi sebelas suku. Peringkatnya:
lebih baik daripada Leibniz dengan selisih yang eksponensial, tetapi di belakang Machin
(yang titik dominannya $\frac15$ lebih kecil daripada $\frac12$):
yaitu kira-kira $0.6$ angka per suku terhadap $1.4$ milik Machin.

\textbf{25.} (i) Untuk $a_k \to 0$ yang menurun, jumlah parsial
berselang-selingnya konvergen dengan $\abs{S - S_n} \leq a_{n+1}$ dan
galatnya membawa tanda suku pertama yang dihilangkan. (ii) Adapun kesamaan hingga
dengan sisa yang eksplisit dapat \emph{dinilai dan
dibatasi di sebuah titik yang dipilih}, sedangkan \omterm{def:b1:logic:statement}{pernyataan} limit hanya
menjanjikan kedekatan yang akhirnya --- sehingga sertifikasi memerlukan yang pertama.
(iii) Yang terpetik: $\pi$ sampai $5\cdot10^{-8}$ lewat Machin, $\ln 2$ sampai
sepuluh desimal lewat deret $\frac13$-nya, keirasionalan $\cos 1$,
$\sin 1$, $\cosh 1$, $\cos\frac1m$ dan $\eu^{-1}$, beserta
kesamaan Taylor--Lagrange, dan peringatan fungsi yang rata. (iv) Adapun
\cref{ch:b1:series} menaikkan pangkat pertanyaan 1 menjadi uji deret
berselang-seling, memisahkan kekonvergenan mutlak dari yang bersyarat,
lalu soal akhir pekannya mementaskan drama penyusunan ulang yang
bagi drama itu deret harmonik berselang-seling Cara 1 menjadi saksi
utamanya.

\end{solution}
